\documentclass[9pt,a4paper]{extarticle}
\usepackage[a4paper,top=1.4cm,bottom=1.6cm,left=1.4cm,right=1.4cm]{geometry}
\usepackage{fontspec}
\setmainfont{Latin Modern Roman}
\setmonofont{DejaVu Sans Mono}[Scale=MatchLowercase]
\usepackage{amsmath,amssymb}
\usepackage{xcolor}
\usepackage{multicol}
\usepackage{array,booktabs}
\usepackage{enumitem}
\usepackage{etoolbox}
\usepackage{needspace}
\usepackage{titlesec}
\usepackage{listings}
\usepackage{microtype}
\usepackage[hidelinks,bookmarksnumbered]{hyperref}
\hypersetup{pdftitle={DSA Notes in Python},pdfauthor={}}

\setlength{\columnsep}{0.7cm}
\setlength{\columnseprule}{0.3pt}
\setlength{\parindent}{0pt}
\setlength{\parskip}{3pt plus 1pt}
\emergencystretch=2em
\clubpenalty=5000 \widowpenalty=5000
\newlength{\tw}
\newsavebox{\tblbox}

\setlist[itemize]{nosep,leftmargin=1.1em,itemsep=1pt,topsep=2pt,partopsep=0pt}
\setlist[itemize,2]{leftmargin=1.2em,label=--}

\setcounter{secnumdepth}{2}
\setcounter{tocdepth}{2}
\titleformat{\section}{\large\bfseries}{\thesection}{0.5em}{}[{\titlerule[0.5pt]}]
\titleformat{\subsection}{\normalsize\bfseries}{\thesubsection}{0.5em}{}
\titleformat{\subsubsection}{\small\bfseries\itshape}{}{0pt}{}
\titlespacing*{\section}{0pt}{9pt plus 2pt}{4pt}
\titlespacing*{\subsection}{0pt}{6pt plus 2pt}{2pt}
\titlespacing*{\subsubsection}{0pt}{4pt plus 1pt}{1pt}
\pretocmd{\section}{\needspace{13\baselineskip}}{}{}
\pretocmd{\subsection}{\needspace{7\baselineskip}}{}{}
\pretocmd{\subsubsection}{\needspace{6\baselineskip}}{}{}

\definecolor{cmt}{gray}{0.38}
\lstset{
  language=Python,
  basicstyle=\ttfamily\fontsize{7}{8.5}\selectfont,
  keywordstyle=\bfseries,
  commentstyle=\itshape\color{cmt},
  morekeywords={nonlocal,True,False,None},
  showstringspaces=false,
  columns=fixed,keepspaces=true,
  breaklines=true,breakindent=1.2em,
  postbreak=\mbox{\textcolor{cmt}{$\hookrightarrow$}\space},
  upquote=true,
  tabsize=4,
  backgroundcolor=\color{black!5},
  frame=none,
  xleftmargin=2pt,xrightmargin=2pt,
  aboveskip=3pt,belowskip=3pt,
}

\begin{document}
\begin{center}
{\LARGE\bfseries DSA Notes in Python}
\end{center}
\vspace{-4pt}
\begin{multicols}{2}
\small
\tableofcontents
\end{multicols}
\newpage
\begin{multicols}{2}
\section{Hash Tables: Dictionaries and Sets}

\subsection{Dictionary}

\begin{lstlisting}
d = {'a': 1, 'b': 2}
\end{lstlisting}

\sbox{\tblbox}{\begingroup\footnotesize\renewcommand{\arraystretch}{1.15}%
\setlength{\tabcolsep}{2.5pt}%
\setlength{\tw}{\dimexpr\columnwidth-6\tabcolsep\relax}%
\begin{tabular}{@{}>{\raggedright\arraybackslash}p{0.2308\tw}>{\raggedright\arraybackslash}p{0.3061\tw}>{\raggedright\arraybackslash}p{0.3543\tw}>{\raggedright\arraybackslash}p{0.1089\tw}@{}}
\toprule
\textbf{Operation} & \textbf{Syntax} & \textbf{Average case} & \textbf{Worst case} \\
\midrule
Access value & \texttt{d[\allowbreak{}'a']} or \texttt{d.\allowbreak{}get(\allowbreak{}'a')} & $O(1)$ & $O(n)$ \\
Insert / update & \texttt{d[\allowbreak{}'c'] =\allowbreak{} 3} & $O(1)$ & $O(n)$ \\
Delete & \texttt{del d[\allowbreak{}'a']} or \texttt{d.\allowbreak{}pop(\allowbreak{}'a')} & $O(1)$ & $O(n)$ \\
Check key exists & \texttt{'a' in d} & $O(1)$ & $O(n)$ \\
Get with default & \texttt{d.\allowbreak{}get(\allowbreak{}'x',\allowbreak{} default)} & $O(1)$ & $O(n)$ \\
Get all keys & \texttt{d.\allowbreak{}keys(\allowbreak{})} & $O(1)$ to get view, $O(n)$ to iterate & — \\
Get all values & \texttt{d.\allowbreak{}values(\allowbreak{})} & $O(1)$ to get view, $O(n)$ to iterate & — \\
Get all items & \texttt{d.\allowbreak{}items(\allowbreak{})} & $O(1)$ to get view, $O(n)$ to iterate & — \\
Length & \texttt{len(\allowbreak{}d)} & $O(1)$ & — \\
Pop last added item & \texttt{d.\allowbreak{}popitem(\allowbreak{})} & $O(1)$ & — \\
Merge & \texttt{d.\allowbreak{}update(\allowbreak{}other)} & $O(\mathrm{len}(other))$ & — \\
Copy & \texttt{d.\allowbreak{}copy(\allowbreak{})} & $O(n)$ & — \\
Clear & \texttt{d.\allowbreak{}clear(\allowbreak{})} & $O(n)$ & — \\
setdefault & \texttt{d.\allowbreak{}setdefault(\allowbreak{}'x',\allowbreak{} [\allowbreak{}])} & $O(1)$ & $O(n)$ \\
\bottomrule
\end{tabular}\endgroup}

\needspace{\dimexpr\ht\tblbox+\dp\tblbox+4.1\baselineskip\relax}

\noindent\textbf{Common Operations}\par\nopagebreak\vspace{1pt}

\usebox{\tblbox}\smallskip

The $O(n)$ worst case only happens with pathological hash collisions; in practice treat these operations as $O(1)$.


\noindent\textbf{Frequency counting}\par\nopagebreak\vspace{1pt}

\begin{lstlisting}
# vanilla dictionary
freq = {}
for x in arr:
    freq[x] = freq.get(x,0) + 1

# with defaultdict
from collections import defaultdict
freq = defaultdict(int) # int as factory means missing keys default to int() == 0
for x in arr:
    freq[x] += 1

# counter (especially for frequency calc)
from collections import Counter
freq = Counter(arr)
\end{lstlisting}

\begin{lstlisting}
from collections import Counter

c = Counter("aabbbc")
print(c) # Counter({'b': 3, 'a': 2, 'c': 1})
print(c.most_common(2)) # [('b', 3), ('a', 2)]
\end{lstlisting}

\subsection{Sets}

\begin{lstlisting}
s = {1, 2, 3}   # initialisation with values
s = set()       # empty set; note that {} creates an empty dict, not a set
\end{lstlisting}

\sbox{\tblbox}{\begingroup\footnotesize\renewcommand{\arraystretch}{1.15}%
\setlength{\tabcolsep}{2.5pt}%
\setlength{\tw}{\dimexpr\columnwidth-6\tabcolsep\relax}%
\begin{tabular}{@{}>{\raggedright\arraybackslash}p{0.2513\tw}>{\raggedright\arraybackslash}p{0.2610\tw}>{\raggedright\arraybackslash}p{0.3788\tw}>{\raggedright\arraybackslash}p{0.1089\tw}@{}}
\toprule
\textbf{Operation} & \textbf{Syntax} & \textbf{Average case} & \textbf{Worst case} \\
\midrule
Add element & \texttt{s.\allowbreak{}add(\allowbreak{}4)} & $O(1)$ & $O(n)$ \\
Remove element & \texttt{s.\allowbreak{}remove(\allowbreak{}4)} (raises \texttt{KeyError} if missing) & $O(1)$ & $O(n)$ \\
Discard element & \texttt{s.\allowbreak{}discard(\allowbreak{}4)} (no error if missing) & $O(1)$ & $O(n)$ \\
Check membership & \texttt{4 in s} & $O(1)$ & $O(n)$ \\
Length & \texttt{len(\allowbreak{}s)} & $O(1)$ & — \\
Pop an arbitrary element & \texttt{s.\allowbreak{}pop(\allowbreak{})} & $O(1)$ & — \\
Union & \texttt{s1 | s2} & $O(\mathrm{len}(s1) + \mathrm{len}(s2))$ & — \\
Intersection & \texttt{s1 \& s2} & $O(\mathrm{min}(\mathrm{len}(s1),\allowbreak\ \mathrm{len}(s2)))$ & — \\
Difference & \texttt{s1 - s2} & $O(\mathrm{len}(s1))$ & — \\
Symmetric difference & \texttt{s1 \^{} s2} & $O(\mathrm{len}(s1) + \mathrm{len}(s2))$ & — \\
Subset check & \texttt{s1 <=\allowbreak{} s2} & $O(\mathrm{len}(s1))$ & — \\
Copy & \texttt{s.\allowbreak{}copy(\allowbreak{})} & $O(n)$ & — \\
\bottomrule
\end{tabular}\endgroup}

\needspace{\dimexpr\ht\tblbox+\dp\tblbox+4.1\baselineskip\relax}

\noindent\textbf{Common Operations}\par\nopagebreak\vspace{1pt}

\usebox{\tblbox}\smallskip

Dictionary keys and set elements must be \emph{hashable}: numbers, strings and tuples (of hashables) are; lists, sets and dicts are not.


\section{Arrays}

\subsection{Static Arrays}

\begin{itemize}
\item Fixed size, decided at the time of creation, allocated as a contiguous block of memory.
\item Access by index is $O(1)$.
\item Insertion/deletion in the middle is $O(n)$, as a lot of values need to be shifted.
\item \texttt{numpy.\allowbreak{}ndarray}s behave like static arrays: resizing means allocating a new array and copying, not growing the old one.
\item Python \texttt{list}s are dynamic arrays (next section).
\end{itemize}

\subsection{Dynamic Arrays}

\begin{itemize}
\item Same contiguous memory idea, but it can resize.
\item If we exceed the capacity, it allocates a bigger block and copies everything over. The textbook choice is to double the size; CPython actually over-allocates by only about 12.5\%, which is still enough for amortized $O(1)$ appends.
\item Append is $O(1)$ \emph{amortized}, not worst-case: the single append that triggers a resize costs $O(n)$, but this averages out over many appends. (A static array has no append at all.)
\item Inserting at the front is $O(n)$ in both cases, as every element has to be shifted.
\end{itemize}

\sbox{\tblbox}{\begingroup\footnotesize\renewcommand{\arraystretch}{1.15}%
\setlength{\tabcolsep}{2.5pt}%
\setlength{\tw}{\dimexpr\columnwidth-6\tabcolsep\relax}%
\begin{tabular}{@{}>{\raggedright\arraybackslash}p{0.1926\tw}>{\raggedright\arraybackslash}p{0.2508\tw}>{\raggedright\arraybackslash}p{0.2113\tw}>{\raggedright\arraybackslash}p{0.3454\tw}@{}}
\toprule
\textbf{Category} & \textbf{Operation} & \textbf{Time Complexity} & \textbf{Notes} \\
\midrule
\textbf{Access} & \texttt{arr[\allowbreak{}i]} & $O(1)$ & Direct index access \\
 & \texttt{len(\allowbreak{}arr)} & $O(1)$ & Python lists track length, not counted each time \\
 & \texttt{arr[\allowbreak{}a:b]} & $O(k)$ & Slicing creates a new list, doesn't view the original \\
\textbf{Adding Elements} & \texttt{arr.\allowbreak{}append(\allowbreak{}x)} & $O(1)$ amortized & Adds at the end \\
 & \texttt{arr.\allowbreak{}extend(\allowbreak{}iterable)} & $O(k)$ amortized & Adds each element of another iterable \\
 & \texttt{arr.\allowbreak{}insert(\allowbreak{}i,\allowbreak{} x)} & $O(n)$ & Everything from the index onwards needs to be shifted \\
\textbf{Removing Elements} & \texttt{arr.\allowbreak{}pop(\allowbreak{})} & $O(1)$ & Removes from the end, no shift needed \\
 & \texttt{arr.\allowbreak{}pop(\allowbreak{}i)} & $O(n)$ & Everything past the index needs to be shifted \\
 & \texttt{arr.\allowbreak{}remove(\allowbreak{}x)} & $O(n)$ & Searches for the first match, then shifts \\
 & \texttt{del arr[\allowbreak{}i]} & $O(n)$ & Same cost as pop \\
 & \texttt{arr.\allowbreak{}clear(\allowbreak{})} & $O(n)$ & — \\
\bottomrule
\end{tabular}\endgroup}

\needspace{\dimexpr\ht\tblbox+\dp\tblbox+4.1\baselineskip\relax}

\subsection{Common Operations on Python Lists}

\usebox{\tblbox}\smallskip

\sbox{\tblbox}{\begingroup\footnotesize\renewcommand{\arraystretch}{1.15}%
\setlength{\tabcolsep}{2.5pt}%
\setlength{\tw}{\dimexpr\columnwidth-6\tabcolsep\relax}%
\begin{tabular}{@{}>{\raggedright\arraybackslash}p{0.1926\tw}>{\raggedright\arraybackslash}p{0.2508\tw}>{\raggedright\arraybackslash}p{0.2113\tw}>{\raggedright\arraybackslash}p{0.3454\tw}@{}}
\toprule
\textbf{Category} & \textbf{Operation} & \textbf{Time Complexity} & \textbf{Notes} \\
\midrule
\textbf{Search} & \texttt{x in arr} & $O(n)$ & Linear search \\
 & \texttt{arr.\allowbreak{}index(\allowbreak{}x)} & $O(n)$ & Returns position \\
 & \texttt{arr.\allowbreak{}count(\allowbreak{}x)} & $O(n)$ & — \\
\textbf{Reordering} & \texttt{arr.\allowbreak{}sort(\allowbreak{})} & $O(n \log n)$ & TimSort under the hood, stable \\
 & \texttt{sorted(\allowbreak{}arr)} & $O(n \log n)$ & Returns a new list, instead of in-place sorting \\
 & \texttt{arr.\allowbreak{}reverse(\allowbreak{})} & $O(n)$ & — \\
 & \texttt{arr[\allowbreak{}::-1]} & $O(n)$ & Same idea, but returns a new list \\
\textbf{Combining} & \texttt{arr1 +\allowbreak{} arr2} & $O(n+m)$ & Creates a new list \\
 & \texttt{arr * k} & $O(nk)$ & — \\
\bottomrule
\end{tabular}\endgroup}

\needspace{\dimexpr\ht\tblbox+\dp\tblbox+1\baselineskip\relax}

\usebox{\tblbox}\smallskip

\textbf{Note}: \texttt{pop(\allowbreak{}0)} and \texttt{insert(\allowbreak{}0,\allowbreak{} x)} are $O(n)$, not $O(1)$, since the whole list has to shift.


\subsection{Implementation of a Dynamic Array}

\begin{lstlisting}
class DynamicArray:

    def __init__(self, capacity: int):
        self.capacity = max(1, capacity)   # a capacity of 0 could never grow by doubling
        self.length = 0
        self.arr = [0]*self.capacity

    def get(self, i: int) -> int:
        if not 0 <= i <= self.length-1:
            raise IndexError('Index out of range')
        return self.arr[i]

    def set(self, i: int, n: int) -> None:
        if not 0 <= i < self.length:
            raise IndexError('Index out of range')
        self.arr[i] = n

    def pushback(self, n: int) -> None:
        if self.length == self.capacity:
            self.resize()

        self.arr[self.length] = n
        self.length += 1

    def popback(self) -> int:
        if self.length == 0:
            raise IndexError('pop from empty array')
        self.length -= 1
        return self.arr[self.length]

    def resize(self) -> None:
        self.capacity = 2*self.capacity
        new_arr = [0]* self.capacity

        for i in range(self.length):
            new_arr[i] = self.arr[i]
        self.arr = new_arr

    def getSize(self) -> int:
        return self.length

    def getCapacity(self) -> int:
        return self.capacity
\end{lstlisting}

\section{Strings}

\subsection{Basics of Strings}

\begin{itemize}
\item Strings are immutable, so every modification on them creates a new string.
\item So, \texttt{s +\allowbreak{}=\allowbreak{} char} in a loop is $O(n^2)$ in the worst case. (CPython sometimes optimises this in place, but that is an implementation detail you should not rely on.)
\end{itemize}

\begin{lstlisting}
s = "hello"
s[0] = 'H' # TypeError: 'str' object does not support item assignment
\end{lstlisting}

\noindent\textbf{Fix for building strings in a loop}\par\nopagebreak\vspace{1pt}

\begin{lstlisting}
# bad method - O(n^2)
result = ""
for c in some_list:
    result += c

# good method - O(n)
result = []
for c in some_list:
    result.append(c)
result = "".join(result)
\end{lstlisting}

\noindent\textbf{Access and Slicing}\par\nopagebreak\vspace{1pt}

\begin{lstlisting}
s = "hello world"

s[0]        # 'h', O(1)
s[-1]       # 'd', O(1)
s[2:5]      # 'llo', O(k), slicing creates a new string
s[::-1]     # reverses the string, O(n)
s[::2]      # 'hlowrd', every second character
len(s)      # length, O(1)
\end{lstlisting}

\subsection{Common Built-in Methods}

\noindent\textbf{Case and Formatting}\par\nopagebreak\vspace{1pt}

\begin{lstlisting}
s = "hello world"
s.upper()         # "HELLO WORLD"
s.lower()         # "hello world"
s.title()         # "Hello World"
s.capitalize()    # "Hello world"
s.swapcase()      # "HELLO WORLD", swaps upper/lower
\end{lstlisting}

\noindent\textbf{Whitespace and cleaning}\par\nopagebreak\vspace{1pt}

\begin{lstlisting}
s.strip()           # removes leading and trailing whitespaces
s.strip('xyz')      # removes leading and trailing chars from the given set
s.lstrip()          # removes leading whitespaces
s.rstrip()          # removes trailing whitespaces
\end{lstlisting}

\noindent\textbf{Searching}\par\nopagebreak\vspace{1pt}

\begin{lstlisting}
s = "hello world"
s.find('lo')         # 3, first occurrence, returns -1 if not found
s.rfind('lo')        # 3, last occurrence, returns -1 if not found
s.index('lo')        # same as find, but raises ValueError if not found
'lo' in s            # O(n), membership check
s.count('l')         # 3, counts non-overlapping occurrences, O(n)
s.startswith('hel')  # O(k)
s.endswith('rld')    # O(k)
\end{lstlisting}

\noindent\textbf{Splitting and Joining}\par\nopagebreak\vspace{1pt}

\begin{lstlisting}
s = 'hello world'

s.split()                   # ['hello', 'world'], splits on white space
s.split(',')                # splits on a specified delimiter
s.split(',', 1)             # maxsplit, only splits once
s.rsplit(',', 1)            # splits from the right
s.splitlines()              # splits on line breaks
','.join(['a', 'b', 'c'])   # 'a,b,c', O(n)
\end{lstlisting}

\texttt{s.\allowbreak{}split(\allowbreak{})} also handles multiple spaces, leading/trailing whitespaces pretty cleanly.


\noindent\textbf{Replacing}\par\nopagebreak\vspace{1pt}

\begin{lstlisting}
s.replace('l', 'L')       # replaces all occurrences, O(n)
s.replace('l', 'L', 1)    # replace only the first occurrence
\end{lstlisting}

\noindent\textbf{Checking character types}\par\nopagebreak\vspace{1pt}

\begin{lstlisting}
c.isalpha()         # is alphabet
c.isdigit()         # is digit
c.isalnum()         # is digit or alphabet
c.isupper()         # is uppercase
c.islower()         # is lowercase
c.isspace()         # is whitespace
\end{lstlisting}

These return \texttt{True} only if \emph{every} character satisfies the condition (and the string is non-empty).


\noindent\textbf{Padding and Alignment}\par\nopagebreak\vspace{1pt}

\begin{lstlisting}
s.zfill(5)          # '00042' for '42', pads with zeros on the left
s.ljust(10, '*')    # left-justify, pad right
s.rjust(10, '*')    # right-justify, pad left
s.center(10, '*')   # center with padding
\end{lstlisting}

\subsection{Character \texorpdfstring{$\leftrightarrow$}{$\leftrightarrow$} ASCII Conversion}

\begin{lstlisting}
ord('a')   # 97, character to ASCII/Unicode code point
chr(97)    # 'a', code point to character

# common pattern: map 'a'-'z' to indices 0-25
ch = 'c'
index = ord(ch) - ord('a')
letter = chr(index + ord('a'))
\end{lstlisting}

\subsection{String \texorpdfstring{$\leftrightarrow$}{$\leftrightarrow$} list conversion (in-place-style manipulation)}

\begin{lstlisting}
s = 'hello'
chars = list(s)     # ['h', 'e', 'l', 'l', 'o']
chars[0] = 'H'
s = ''.join(chars)  # 'Hello'
\end{lstlisting}

\subsection{String Comparison}

\begin{lstlisting}
"abc" == "abc"          # True, value comparison
"abc" < "abd"           # True, lexicographic comparison, character by character

sorted("cab")           # ['a', 'b', 'c'], returns a list, not a string
''.join(sorted("cab"))  # 'abc'
\end{lstlisting}

\end{multicols}

\sbox{\tblbox}{\begingroup\footnotesize\renewcommand{\arraystretch}{1.15}%
\setlength{\tabcolsep}{2.5pt}%
\setlength{\tw}{\dimexpr\textwidth-10\tabcolsep\relax}%
\begin{tabular}{@{}>{\raggedright\arraybackslash}p{0.2915\tw}>{\raggedright\arraybackslash}p{0.1178\tw}>{\raggedright\arraybackslash}p{0.1234\tw}>{\raggedright\arraybackslash}p{0.1193\tw}>{\raggedright\arraybackslash}p{0.0935\tw}>{\raggedright\arraybackslash}p{0.2546\tw}@{}}
\toprule
\textbf{Algorithm} & \textbf{Best} & \textbf{Average} & \textbf{Worst} & \textbf{Space} & \textbf{Stable?} \\
\midrule
Bubble Sort & $O(n)$ & $O(n^2)$ & $O(n^2)$ & $O(1)$ & Yes \\
Selection Sort & $O(n^2)$ & $O(n^2)$ & $O(n^2)$ & $O(1)$ & No \\
Insertion Sort & $O(n)$ & $O(n^2)$ & $O(n^2)$ & $O(1)$ & Yes \\
Merge Sort & $O(n \log n)$ & $O(n \log n)$ & $O(n \log n)$ & $O(n)$ & Yes \\
Quick Sort & $O(n \log n)$ & $O(n \log n)$ & $O(n^2)$ & $O(\log n)$ & No \\
Heap Sort & $O(n \log n)$ & $O(n \log n)$ & $O(n \log n)$ & $O(1)$ & No \\
Counting Sort & $O(n+k)$ & $O(n+k)$ & $O(n+k)$ & $O(n+k)$ & Yes \\
Radix Sort & $O(nk)$ & $O(nk)$ & $O(nk)$ & $O(n+k)$ & Yes \\
Bucket Sort & $O(n+k)$ & $O(n+k)$ & $O(n^2)$ & $O(n)$ & Yes (depends on inner sort) \\
Tim Sort (Python's \texttt{sort(\allowbreak{})}) & $O(n)$ & $O(n \log n)$ & $O(n \log n)$ & $O(n)$ & Yes \\
\bottomrule
\end{tabular}\endgroup}

\needspace{\dimexpr\ht\tblbox+\dp\tblbox+6.7\baselineskip\relax}

\section{Sorting}

\subsection{Time Complexities}

\usebox{\tblbox}\smallskip

\begin{multicols}{2}

\textbf{Stability:} a sort is stable if elements with equal keys keep their original relative order (this matters when sorting records, e.g. by age after sorting by name).


Notes on the table: the quick sort space is the average recursion depth (worst case $O(n)$); bubble and insertion sort reach $O(n)$ only on (nearly) sorted input; counting sort is stable only in its prefix-sum form (shown under radix sort); $k$ is the value range for counting/bucket sort and the number of digits for radix sort.


\subsection{Comparison based Algorithms}

\subsubsection{Bubble Sort}

\begin{itemize}
\item Repeatedly swaps adjacent elements if they're out of order.
\item Each pass “bubbles” the largest unsorted element to its correct position at the end.
\end{itemize}

\begin{lstlisting}
def bubble_sort(arr):
    n = len(arr)
    for i in range(n):
        swapped = False
        for j in range(n-i-1):
            if arr[j] > arr[j+1]:
                arr[j], arr[j+1] = arr[j+1], arr[j]
                swapped = True
        if not swapped:
            break  # already sorted, stop early
    return arr
\end{lstlisting}

\subsubsection{Selection Sort}

\begin{itemize}
\item Repeatedly finds the minimum of the unsorted portion and swaps it into place.
\end{itemize}

\begin{lstlisting}
def selection_sort(arr):
    n = len(arr)
    for i in range(n):
        min_idx = i
        for j in range(i+1,n):
            if arr[j] < arr[min_idx]:
                min_idx = j
        arr[i], arr[min_idx] = arr[min_idx], arr[i]
    return arr
\end{lstlisting}

\subsubsection{Insertion Sort}

\begin{itemize}
\item Builds the sorted array, one element at a time.
\item Inserts each new element into its correct position among the already-sorted part.
\end{itemize}

\begin{lstlisting}
def insertion_sort(arr):
    for i in range(1,len(arr)):
        key = arr[i]
        j = i-1
        while j>= 0 and arr[j] > key:
            arr[j+1] = arr[j]
            j -= 1
        arr[j+1] = key
    return arr
\end{lstlisting}

\subsubsection{Merge Sort}

\begin{itemize}
\item DnC: Divide the array into halves, sort them and merge back.
\item Stable (because of the \texttt{<=\allowbreak{}} in the merge step), but needs $O(n)$ extra space.
\end{itemize}

\begin{lstlisting}
def merge_sort(arr):
    if len(arr) <= 1:
        return arr

    mid = len(arr)//2

    left = merge_sort(arr[:mid])
    right = merge_sort(arr[mid:])

    return merge(left,right)

def merge(left, right):
    result = []
    i = j = 0

    while i < len(left) and j < len(right):
        if left[i] <= right[j]:
            result.append(left[i])
            i += 1
        else:
            result.append(right[j])
            j += 1

    result.extend(left[i:])
    result.extend(right[j:])

    return result
\end{lstlisting}

\subsubsection{Quick Sort}

\begin{itemize}
\item DnC: Pick a pivot, partition the array such that smaller elements are to the left of the pivot and the rest are on the right.
\item Recurse on each step.
\item With the last element as pivot, already sorted input is the worst case ($O(n^2)$ time and $O(n)$ recursion depth); picking a random pivot avoids this in practice.
\end{itemize}

\begin{lstlisting}
def quick_sort(arr, low=0, high=None):
    if high is None:
        high = len(arr)-1
    if low < high:
        pivot_idx = partition(arr,low,high)
        quick_sort(arr,low,pivot_idx-1)
        quick_sort(arr,pivot_idx+1,high)
    return arr

def partition(arr, low, high):
    pivot = arr[high]
    i = low - 1
    for j in range(low, high):
        if arr[j] <= pivot:
            i += 1
            arr[i], arr[j] = arr[j], arr[i]
    arr[i+1], arr[high] = arr[high], arr[i+1]
    return i+1
\end{lstlisting}

\subsubsection{Heap Sort}

\begin{itemize}
\item Builds a max-heap from the array, then repeatedly extract the max and place it at the end.
\end{itemize}

\begin{lstlisting}
def heap_sort(arr):
    n = len(arr)
    for i in range(n//2 - 1, -1, -1):
        heapify(arr, n, i)
    for i in range(n-1, 0, -1):
        arr[0], arr[i] = arr[i], arr[0]
        heapify(arr, i, 0)
    return arr

def heapify(arr, n, i):
    largest = i
    left, right = 2*i + 1, 2*i + 2

    if left < n and arr[left] > arr[largest]:
        largest = left
    if right < n and arr[right] > arr[largest]:
        largest = right
    if largest != i:
        arr[i], arr[largest] = arr[largest], arr[i]
        heapify(arr, n, largest)
\end{lstlisting}

\subsection{Non-comparison based Algorithms}

\subsubsection{Counting Sort}

\begin{itemize}
\item Counts the occurrences of each value and then reconstructs the sorted array. Works for non-negative integers; $k$ is the maximum value.
\end{itemize}

\begin{lstlisting}
def counting_sort(arr):
    if not arr: return arr

    max_val = max(arr)
    count = [0]*(max_val + 1)

    for x in arr:
        count[x] += 1

    result = []
    for val, c in enumerate(count):
        result.extend([val]*c)
    return result
\end{lstlisting}

\subsubsection{Radix Sort}

\begin{itemize}
\item Sorts integers digit by digit, least significant digit first, using a \textbf{stable} counting sort on each digit. $O(d \cdot (n + b))$ for $d$ digits in base $b$.
\item The prefix-sum form of counting sort below is what makes counting sort stable: \texttt{count[\allowbreak{}d]} becomes the end position of digit \texttt{d}, and the output is filled from the back.
\end{itemize}

\begin{lstlisting}
def radix_sort(arr):
    """Non-negative integers only."""
    if not arr: return arr
    max_val = max(arr)
    exp = 1
    while max_val // exp > 0:
        arr = counting_sort_by_digit(arr, exp)
        exp *= 10
    return arr

def counting_sort_by_digit(arr, exp):
    count = [0]*10
    for x in arr:
        count[(x // exp) % 10] += 1
    for d in range(1, 10):
        count[d] += count[d-1]           # prefix sums
    out = [0]*len(arr)
    for x in reversed(arr):              # go backwards to stay stable
        d = (x // exp) % 10
        count[d] -= 1
        out[count[d]] = x
    return out
\end{lstlisting}

\subsubsection{Bucket Sort}

\begin{itemize}
\item Spread the values into buckets by range, sort every bucket, then concatenate. $O(n+k)$ on average for roughly uniformly distributed input, $O(n^2)$ if everything lands in one bucket.
\end{itemize}

\begin{lstlisting}
def bucket_sort(arr):
    """Floats in [0, 1)."""
    n = len(arr)
    if n == 0: return arr
    buckets = [[] for _ in range(n)]
    for x in arr:
        buckets[int(x * n)].append(x)
    result = []
    for b in buckets:
        result.extend(sorted(b))
    return result
\end{lstlisting}

\section{Linked Lists}

\subsection{Basic Structure}

\begin{itemize}
\item Nodes hold a value and a pointer to the next node; there is no contiguous memory and no random access.
\item \texttt{get(\allowbreak{}i)}, \texttt{insertTail} and \texttt{remove(\allowbreak{}i)} are $O(n)$; \texttt{insertHead} is $O(1)$ (keeping a tail pointer would make \texttt{insertTail} $O(1)$ too).
\end{itemize}

\begin{lstlisting}
# linked lists
from typing import List

class Node:
    def __init__(self, val=None):
        self.val = val
        self.next = None

class LinkedList:
    def __init__(self):
        self.head = None

    def get(self, index: int) -> int:
        if index < 0:
            return -1
        temp = self.head
        i = 0
        while i < index and temp:
            temp = temp.next
            i += 1
        return temp.val if temp else -1

    def insertHead(self, val: int) -> None:
        temp = self.head
        self.head = Node(val)
        self.head.next = temp

    def insertTail(self, val: int) -> None:
        if not self.head:
            self.head = Node(val)
            return
        temp = self.head
        while temp.next:
            temp = temp.next
        temp.next = Node(val)

    def remove(self, index: int) -> bool:
        if not self.head or index < 0:
            return False
        if index == 0:
            self.head = self.head.next
            return True
        curr = self.head
        i = 0
        while curr.next and i < index - 1:
            curr = curr.next
            i += 1
        if curr.next is None:
            return False
        curr.next = curr.next.next
        return True

    def getValues(self) -> List[int]:
        result = []
        curr = self.head
        while curr:
            result.append(curr.val)
            curr = curr.next
        return result
\end{lstlisting}

\subsection{Reversing a Linked List}

\begin{lstlisting}
def reverseLL(head: Node) -> Node:
    prev, current = None, head

    while current:
        nextNode = current.next
        current.next = prev
        prev = current
        current = nextNode

    return prev
\end{lstlisting}

\subsection{Detecting Cycle}

\begin{itemize}
\item Floyd's tortoise and hare: \texttt{slow} moves one step, \texttt{fast} moves two; they meet if and only if the list has a cycle. $O(n)$ time, $O(1)$ space.
\end{itemize}

\begin{lstlisting}
def detectCycle(head: Node) -> bool:
    slow = fast = head

    while fast and fast.next:
        slow = slow.next
        fast = fast.next.next

        if slow == fast: return True

    return False
\end{lstlisting}

\section{Queue}

\subsection{Introduction}

\begin{itemize}
\item FIFO (First In, First Out) functionality
\item Heavily used in Graph BFS Traversals. \texttt{list.\allowbreak{}pop(\allowbreak{}0)} is $O(n)$, so a plain list is a poor fit.
\item We mostly use the \texttt{deque} (pronounced as deck) functionality from \texttt{collections} library.
\end{itemize}

\subsection{Some important features/functions on \texorpdfstring{\texttt{deque}}{deque}}

\begin{lstlisting}
from collections import deque
d = deque([1,2,3])
\end{lstlisting}

\sbox{\tblbox}{\begingroup\footnotesize\renewcommand{\arraystretch}{1.15}%
\setlength{\tabcolsep}{2.5pt}%
\setlength{\tw}{\dimexpr\columnwidth-4\tabcolsep\relax}%
\begin{tabular}{@{}>{\raggedright\arraybackslash}p{0.3390\tw}>{\raggedright\arraybackslash}p{0.2068\tw}>{\raggedright\arraybackslash}p{0.4543\tw}@{}}
\toprule
\textbf{Operation} & \textbf{Time Complexity} & \textbf{Notes} \\
\midrule
\texttt{d.\allowbreak{}append(\allowbreak{}4)} & $O(1)$ & Appends to the right \\
\texttt{d.\allowbreak{}appendleft(\allowbreak{}0)} & $O(1)$ & Appends to the left \\
\texttt{d.\allowbreak{}pop(\allowbreak{})} & $O(1)$ & Removes and returns the rightmost element \\
\texttt{d.\allowbreak{}popleft(\allowbreak{})} & $O(1)$ & Removes and returns the leftmost element \\
\texttt{d.\allowbreak{}extend(\allowbreak{}iterable)} & $O(k)$ & Appends multiple elements to the right \\
\texttt{d.\allowbreak{}extendleft(\allowbreak{}iterable)} & $O(k)$ & Appends multiple elements to the left (in reverse order) \\
\texttt{d.\allowbreak{}insert(\allowbreak{}i,\allowbreak{} x)} & $O(n)$ & Inserts x at index i \\
\texttt{d.\allowbreak{}remove(\allowbreak{}x)} & $O(n)$ & Removes the first occurrence of x \\
\bottomrule
\end{tabular}\endgroup}

\needspace{\dimexpr\ht\tblbox+\dp\tblbox+1.5\baselineskip\relax}

\usebox{\tblbox}\smallskip

\sbox{\tblbox}{\begingroup\footnotesize\renewcommand{\arraystretch}{1.15}%
\setlength{\tabcolsep}{2.5pt}%
\setlength{\tw}{\dimexpr\columnwidth-4\tabcolsep\relax}%
\begin{tabular}{@{}>{\raggedright\arraybackslash}p{0.3390\tw}>{\raggedright\arraybackslash}p{0.2068\tw}>{\raggedright\arraybackslash}p{0.4543\tw}@{}}
\toprule
\textbf{Operation} & \textbf{Time Complexity} & \textbf{Notes} \\
\midrule
\texttt{d.\allowbreak{}rotate(\allowbreak{}n)} & $O(k)$ & Rotates by n steps (negative n means rotate left) \\
\texttt{d.\allowbreak{}reverse(\allowbreak{})} & $O(n)$ & Reverses the deque in place \\
\texttt{d.\allowbreak{}clear(\allowbreak{})} & $O(n)$ & Removes all elements \\
\texttt{d.\allowbreak{}count(\allowbreak{}x)} & $O(n)$ & Counts all occurrences of x \\
\texttt{d.\allowbreak{}index(\allowbreak{}x)} & $O(n)$ & Finds the index of x \\
\texttt{d.\allowbreak{}copy(\allowbreak{})} & $O(n)$ & Shallow copy of the deque \\
\texttt{d[\allowbreak{}i]} & $O(n)$ & Access by index, deque is not optimized for random access \\
\texttt{len(\allowbreak{}d)} & $O(1)$ & Length of the deque \\
\bottomrule
\end{tabular}\endgroup}

\needspace{\dimexpr\ht\tblbox+\dp\tblbox+1\baselineskip\relax}

\usebox{\tblbox}\smallskip

\subsection{An important note}

\begin{itemize}
\item \texttt{deque(\allowbreak{}maxlen=\allowbreak{}N)} creates a bounded \texttt{deque}, which is heavily used in sliding window approaches. If we append beyond \texttt{maxlen}, elements from the opposite end are discarded.
\end{itemize}

\begin{lstlisting}
from collections import deque

d = deque(maxlen=3)
d.append(1); d.append(2); d.append(3); d.append(4)

print(d) # deque([2, 3, 4], maxlen=3), the 1 got dropped off
\end{lstlisting}

\subsection{Implementation of Queues using LinkedLists}

\begin{lstlisting}
class Node:
    def __init__(self, val):
        self.val = val
        self.next = None

class Queue:
    def __init__(self):
        self.front = None
        self.rear = None
        self.size = 0

    def enqueue(self, val):
        newNode = Node(val)
        if not self.rear:
            self.front = self.rear = newNode
        else:
            self.rear.next = newNode
            self.rear = newNode
        self.size += 1

    def dequeue(self) -> int:
        if not self.front:
            raise IndexError("dequeue from an empty queue")
        val = self.front.val
        self.front = self.front.next
        if not self.front:
            self.rear = None
        self.size -= 1
        return val

    def peek(self) -> int:
        if not self.front:
            raise IndexError("peek from an empty queue")
        return self.front.val

    def is_empty(self) -> bool:
        return self.front is None
\end{lstlisting}

\subsection{Implementation of Queues using two stacks}

\begin{itemize}
\item \texttt{dequeue} and \texttt{peek} are $O(1)$ amortized: every element is moved from \texttt{in\_\allowbreak{}stack} to \texttt{out\_\allowbreak{}stack} at most once.
\end{itemize}

\begin{lstlisting}
class QueueUsingStacks:
    def __init__(self):
        self.in_stack = []
        self.out_stack = []

    def enqueue(self, val) -> None:
        self.in_stack.append(val)

    def dequeue(self) -> int:
        if not self.out_stack:
            while self.in_stack:
                self.out_stack.append(self.in_stack.pop())
        if not self.out_stack:
            raise IndexError("dequeue from an empty queue")
        return self.out_stack.pop()

    def peek(self) -> int:
        if not self.out_stack:
            while self.in_stack:
                self.out_stack.append(self.in_stack.pop())
        if not self.out_stack:
            raise IndexError("peek from an empty queue")
        return self.out_stack[-1]

    def is_empty(self) -> bool:
        return not self.in_stack and not self.out_stack
\end{lstlisting}

\subsection{Implementation of Circular Queue}

\begin{itemize}
\item A fixed-capacity ring buffer: \texttt{front} plus \texttt{count} determine the rear index, which avoids the empty-vs-full ambiguity of a two-pointer design.
\end{itemize}

\begin{lstlisting}
class CircularQueue:
    def __init__(self, k):
        self.q = [None]*k
        self.capacity = k
        self.front = 0
        self.count = 0

    def enqueue(self,val) -> None:
        if self.count == self.capacity:
            raise OverflowError("queue is full")
        rear = (self.front + self.count) % self.capacity
        self.q[rear] = val
        self.count += 1

    def dequeue(self) -> int:
        if self.count == 0:
            raise IndexError("Empty Queue")
        val = self.q[self.front]
        self.front = (self.front + 1) % self.capacity
        self.count -= 1
        return val

    def is_empty(self) -> bool:
        return self.count == 0

    def is_full(self) -> bool:
        return self.count == self.capacity
\end{lstlisting}

\section{Double-ended Queues}

\subsection{Introduction}

\begin{itemize}
\item Double ended queues give us $O(1)$ functionality to append/pop on both left and right ends.
\item Python's \texttt{collections.\allowbreak{}deque} is a double ended queue.
\end{itemize}

\subsection{Implementation using Doubly Linked Lists}

\begin{lstlisting}
class Node:
    def __init__(self, val):
        self.val = val
        self.next = None
        self.prev = None

class Deque:
    def __init__(self):
        self.front = None
        self.rear = None
        self.size = 0

    def append(self,val: int) -> None:
        newNode = Node(val)
        if not self.rear:
            self.front = self.rear = newNode
        else:
            newNode.prev = self.rear
            self.rear.next = newNode
            self.rear = newNode
        self.size += 1

    def appendleft(self, val: int) -> None:
        newNode = Node(val)
        if not self.rear:
            self.front = self.rear = newNode
        else:
            newNode.next = self.front
            self.front.prev = newNode
            self.front = newNode
        self.size += 1

    def pop(self) -> int:
        if not self.rear:
            raise IndexError("empty queue")
        val = self.rear.val
        self.rear = self.rear.prev
        if not self.rear:
            self.front = None
        else:
            self.rear.next = None
        self.size -= 1
        return val

    def popleft(self) -> int:
        if not self.rear:
            raise IndexError("empty queue")
        val = self.front.val
        self.front = self.front.next
        if not self.front:
            self.rear = None
        else:
            self.front.prev = None
        self.size -= 1
        return val

    def is_empty(self) -> bool:
        return self.size == 0

# all operations here are in O(1)
\end{lstlisting}

\section{Stack}

\subsection{Stack Basics}

\begin{itemize}
\item Stacks follow the Last In, First Out (LIFO) protocol.
\item We can generally use Python's default Lists as stacks.
\item We may also use \texttt{collections.\allowbreak{}deque}.
\item \texttt{append} and \texttt{pop} at the end of a list are $O(1)$ (amortized for \texttt{append}), so a plain list is enough for most problems.
\end{itemize}

\begin{lstlisting}
stack = []
stack.append(1)
stack.append(2)
stack.append(3)

print(stack.pop()) # pop and return the top element
print(stack[-1]) # peek
print(len(stack)) #size
print(not stack) #is_empty check
\end{lstlisting}

\begin{lstlisting}
# stacks using collections.deque
from collections import deque

stack = deque()
stack.append(1)
stack.append(2)
stack.append(3)

print(stack.pop()) # 3
print(stack[-1]) #peek
\end{lstlisting}

\subsection{Implementation using LinkedLists}

\begin{lstlisting}
class Node:
    def __init__(self, val):
        self.val = val
        self.next = None

class Stack:
    def __init__(self):
        self.top = None
        self.size = 0

    def push(self, val: int) -> None:
        newNode = Node(val)
        newNode.next = self.top
        self.top = newNode
        self.size += 1

    def pop(self) -> int:
        if not self.top:
            raise IndexError("Empty Stack")
        val = self.top.val
        self.top = self.top.next
        self.size -= 1
        return val

    def peek(self) -> int:
        if not self.top:
            raise IndexError("Empty Stack")
        return self.top.val

    def is_empty(self) -> bool:
        return self.top is None
\end{lstlisting}

\section{Heaps / Priority Queues}

\subsection{Basics}

\begin{itemize}
\item A binary heap is a complete binary tree in which every parent is $\le$ its children (min-heap), so the minimum always sits at the root.
\item Unlike LIFO or FIFO, PQs don't care about the insertion order but rather the priority order.
\end{itemize}

\subsection{\texorpdfstring{\texttt{heapq}}{heapq} module}

\begin{itemize}
\item \texttt{heapq} works directly on python \texttt{list}, turning it into a binary heap in-place.
\end{itemize}

\begin{lstlisting}
import heapq

heap = []

heapq.heappush(heap, 5)
heapq.heappush(heap, 1)
heapq.heappush(heap, 3)

print(heapq.heappop(heap)) # 1 -> smallest
print(heapq.heappop(heap)) # 3
\end{lstlisting}

\sbox{\tblbox}{\begingroup\footnotesize\renewcommand{\arraystretch}{1.15}%
\setlength{\tabcolsep}{2.5pt}%
\setlength{\tw}{\dimexpr\columnwidth-4\tabcolsep\relax}%
\begin{tabular}{@{}>{\raggedright\arraybackslash}p{0.3966\tw}>{\raggedright\arraybackslash}p{0.2068\tw}>{\raggedright\arraybackslash}p{0.3967\tw}@{}}
\toprule
\textbf{Function} & \textbf{Time Complexity} & \textbf{Notes} \\
\midrule
\texttt{heapq.\allowbreak{}heappush(\allowbreak{}heap,\allowbreak{} x)} & $O(\log n)$ & Push $x$ onto heap \\
\texttt{heapq.\allowbreak{}heappop(\allowbreak{}heap)} & $O(\log n)$ & Pop and return smallest item \\
\texttt{heapq.\allowbreak{}heappushpop(\allowbreak{}heap,\allowbreak{} x)} & $O(\log n)$ & Push $x$ and then pop the smallest (more efficient than push + pop) \\
\texttt{heapq.\allowbreak{}heapreplace(\allowbreak{}heap,\allowbreak{} x)} & $O(\log n)$ & Pop the smallest, and then push $x$ \\
\texttt{heapq.\allowbreak{}heapify(\allowbreak{}list)} & $O(n)$ & Convert an existing list into a heap, in-place \\
\texttt{heapq.\allowbreak{}nlargest(\allowbreak{}k,\allowbreak{} iterable)} & $O(n \log k)$ & Returns the $k$ largest elements \\
\texttt{heapq.\allowbreak{}nsmallest(\allowbreak{}k,\allowbreak{} iterable)} & $O(n \log k)$ & Returns the smallest $k$ elements \\
\texttt{heap[\allowbreak{}0]} & $O(1)$ & Peek (smallest element sits at index 0, not -1) \\
\bottomrule
\end{tabular}\endgroup}

\needspace{\dimexpr\ht\tblbox+\dp\tblbox+4.1\baselineskip\relax}

\noindent\textbf{heapq functions}\par\nopagebreak\vspace{1pt}

\usebox{\tblbox}\smallskip

\noindent\textbf{max-heap trick}\par\nopagebreak\vspace{1pt}

\begin{itemize}
\item \texttt{heapq} only gives us min-heap, for max-heap, we can simply negate the values on the way in and way out.
\end{itemize}

\begin{lstlisting}
import heapq

heap = []
heapq.heappush(heap, -5)
heapq.heappush(heap, -1)
heapq.heappush(heap, -3)

print(-heapq.heappop(heap)) # 5 -> largest
print(-heapq.heappop(heap)) # 3 -> second largest
\end{lstlisting}

\noindent\textbf{Priority Queues with custom order (tuple)}\par\nopagebreak\vspace{1pt}

\begin{itemize}
\item tuples should be of the format \texttt{(\allowbreak{}priority,\allowbreak{} value)}
\end{itemize}

\begin{lstlisting}
import heapq

heap = []

heapq.heappush(heap, (2, "task x"))
heapq.heappush(heap, (1, "task l"))
heapq.heappush(heap, (3, "task c"))

print(heapq.heappop(heap)) # (1, 'task l')
\end{lstlisting}

\begin{itemize}
\item If two priorities tie, Python compares the next tuple element, which raises \texttt{TypeError} for non-comparable items (e.g. custom objects). Use \texttt{(\allowbreak{}priority,\allowbreak{} counter,\allowbreak{} item)} with an increasing counter as a tie-breaker.
\end{itemize}

\subsection{Implementation of heaps from scratch}

\begin{itemize}
\item A binary heap is stored as an array, where the children of index $i$ element are at $2i + 1$, and $2i + 2$, and its parent lives at $(i-1)//2$.
\item \texttt{push} adds a value at the end and bubbles up (sift up).
\item \texttt{pop} swaps the last value with the root and bubbles down (sift down).
\item Both are $O(\log n)$, as the height of the tree is $\log n$.
\end{itemize}

\begin{lstlisting}
class MinHeap:
    def __init__(self):
        self.heap = []

    def _parent(self, i):
        return (i-1)//2

    def _left(self, i):
        return 2*i+1

    def _right(self, i):
        return 2*i+2

    def push(self, val):
        self.heap.append(val)
        self._sift_up(len(self.heap)-1)

    def _sift_up(self, i):
        while i> 0 and self.heap[i] < self.heap[self._parent(i)]:
            p = self._parent(i)
            self.heap[i], self.heap[p] = self.heap[p], self.heap[i]
            i = p

    def pop(self):
        if not self.heap:
            raise IndexError("Empty Heap")
        top = self.heap[0]
        last = self.heap.pop()
        if self.heap:
            self.heap[0] = last
            self._sift_down(0)
        return top

    def _sift_down(self, i):
        n = len(self.heap)
        while True:
            smallest = i
            l, r = self._left(i), self._right(i)
            if l < n and self.heap[l] < self.heap[smallest]:
                smallest = l
            if r < n and self.heap[r] < self.heap[smallest]:
                smallest = r
            if smallest == i:
                break
            self.heap[i], self.heap[smallest] = self.heap[smallest], self.heap[i]
            i = smallest

    def peek(self):
        if not self.heap:
            raise IndexError("Empty Heap")
        return self.heap[0]

    def is_empty(self):
        return len(self.heap) == 0
\end{lstlisting}

\section{Binary Trees}

\subsection{Binary Trees Basics}

\begin{lstlisting}
# class for TreeNode
class TreeNode:
    def __init__(self, val):
        self.val = val
        self.left = None
        self.right = None

# adding stuff
tree = TreeNode(1)
tree.left = TreeNode(2)
tree.right = TreeNode(3)
tree.left.left = TreeNode(4)
\end{lstlisting}

\subsection{Binary Trees Traversal}

\begin{lstlisting}
# Preorder Traversal - DFS
def preorder_dfs(root: TreeNode) -> None:
    if root:
        print(root.val, end=" ")
        preorder_dfs(root.left)
        preorder_dfs(root.right)
\end{lstlisting}

\begin{lstlisting}
# Inorder Traversal - DFS
def inorder_dfs(root: TreeNode) -> None:
    if root:
        inorder_dfs(root.left)
        print(root.val, end=" ")
        inorder_dfs(root.right)
\end{lstlisting}

\textbf{Note} - Inorder traversal of a BST gives a sorted array.


\begin{lstlisting}
# Postorder Traversal - DFS
def postorder_dfs(root: TreeNode) -> None:
    if root:
        postorder_dfs(root.left)
        postorder_dfs(root.right)
        print(root.val, end=" ")
\end{lstlisting}

\begin{lstlisting}
# BFS Traversal using Queue
from collections import deque
def bfs(root: TreeNode) -> None:
    if not root: return

    queue = deque([root])
    while queue:
        node = queue.popleft()
        print(node.val, end=" ")
        if node.left: queue.append(node.left)
        if node.right: queue.append(node.right)
\end{lstlisting}

\begin{lstlisting}
# iterative inorder traversal using explicit stack
from typing import List

def iterative_inorder(root: TreeNode) -> List[int]:
    result, stack = [], []
    curr = root
    while curr or stack:
        while curr:
            stack.append(curr)
            curr = curr.left
        curr = stack.pop()
        result.append(curr.val)
        curr = curr.right
    return result
\end{lstlisting}

\subsection{Constructing a Binary Tree from preorder and inorder traversals}

\begin{itemize}
\item Preorder gives the root first; its position in the inorder sequence splits the remaining nodes into the left and right subtrees. Assumes all values are unique.
\item The naive version is $O(n^2)$ (\texttt{index} plus slicing); the hashmap version below is $O(n)$.
\end{itemize}

\begin{lstlisting}
# naive implementation
class TreeNode:
    def __init__(self, val):
        self.val = val
        self.left = None
        self.right = None

def build_tree(preorder, inorder):
    if not preorder or not inorder:
        return None

    root_val = preorder[0]
    root = TreeNode(root_val)

    mid = inorder.index(root_val)

    root.left = build_tree(preorder[1:mid+1], inorder[:mid])
    root.right = build_tree(preorder[mid+1:], inorder[mid+1:])

    return root
\end{lstlisting}

\begin{lstlisting}
# optimal implementation
# we use a hashmap to lookup inorder index and used only pointers as opposed to copying the whole arrays

def build_tree(preorder, inorder):
    inorder_index = {val: i for i, val in enumerate(inorder)}
    pre_idx = 0

    def helper(left, right):
        nonlocal pre_idx
        if left > right:
            return None

        root_val = preorder[pre_idx]
        root = TreeNode(root_val)
        pre_idx += 1

        mid = inorder_index[root_val]

        root.left = helper(left, mid-1)
        root.right = helper(mid+1, right)

        return root

    return helper(0, len(preorder)-1)
\end{lstlisting}

\subsection{Height of a Binary Tree}

\begin{itemize}
\item Here the height counts \textbf{nodes} on the longest root-to-leaf path (an empty tree has height 0). To count edges instead, return \texttt{-1} for an empty tree.
\end{itemize}

\begin{lstlisting}
def height(root: TreeNode) -> int:
    if not root: return 0

    leftHt = height(root.left)
    rightHt = height(root.right)

    return 1 + max(leftHt, rightHt)
\end{lstlisting}

\subsection{Number of Nodes in a Binary Tree}

\begin{lstlisting}
def numNodes(root: TreeNode) -> int:
    if not root: return 0

    leftNums = numNodes(root.left)
    rightNums = numNodes(root.right)

    return 1 + leftNums + rightNums
\end{lstlisting}

\subsection{Is Identical}

\begin{lstlisting}
def isIdentical(p: TreeNode, q: TreeNode) -> bool:
    if not p or not q: return p == q

    return (
        p.val ==  q.val and
        isIdentical(p.left, q.left) and
        isIdentical(p.right, q.right)
    )
\end{lstlisting}

\subsection{Is Subtree}

\begin{itemize}
\item Test \texttt{isIdentical} at every node of \texttt{root}: $O(m \cdot n)$ worst case for trees with $n$ and $m$ nodes.
\end{itemize}

\begin{lstlisting}
def isSubtree(root: TreeNode, subRoot: TreeNode) -> bool:

    if not subRoot: return True
    if not root: return False

    if (root.val == subRoot.val and isIdentical(root, subRoot)): return True

    return (
        isSubtree(root.left, subRoot) or
        isSubtree(root.right, subRoot)
    )
\end{lstlisting}

\subsection{Diameter of a binary tree}

\begin{itemize}
\item Maximum distance (in edges) between any two nodes; the path may or may not pass through the root. The longest path through a node is \texttt{left\_\allowbreak{}height +\allowbreak{} right\_\allowbreak{}height}.
\end{itemize}

\begin{lstlisting}
def diameter(root: TreeNode) -> int:
    result = 0
    def depth(root: TreeNode) -> int:
        nonlocal result
        if not root:
            return 0
        left = depth(root.left)
        right = depth(root.right)
        result = max(result, left+right)
        return 1 + max(left, right)
    depth(root)
    return result
\end{lstlisting}

\subsection{Check if balanced}

\begin{itemize}
\item Balanced tree: for every node, the height difference between left and right subtrees is at most 1.
\item Naive approach recomputes height at every node $\to$ $O(n^2)$. Better: compute height and check balance in a single bottom-up pass $\to$ $O(n)$.
\end{itemize}

\begin{lstlisting}
def is_balanced(root: TreeNode) -> bool:
    def height(node: TreeNode) -> int:
        if not node:
            return 0
        left = height(node.left)
        if left == -1:
            return -1
        right = height(node.right)
        if right == -1:
            return -1
        if abs(left - right) > 1:
            return -1
        return 1 + max(left, right)

    return height(root) != -1
\end{lstlisting}

\section{Binary Search Trees}

\subsection{Basics}

\begin{itemize}
\item For every node, all values in the left subtree are smaller and all values in the right subtree are larger (\texttt{left < node < right}). We assume distinct keys; the \texttt{insert} below ignores duplicates.
\item Search, insert and delete are $O(h)$ for height $h$: $O(\log n)$ when balanced, $O(n)$ for a degenerate tree (e.g. keys inserted in sorted order).
\end{itemize}

\subsection{Implementation from Scratch}

\begin{lstlisting}
class TreeNode:
    def __init__(self, val: int):
        self.val = val
        self.left = None
        self.right = None

class BST:
    def __init__(self):
        self.root = None

    def insert(self, val: int) -> None:
        self.root = self._insert(self.root, val)

    def _insert(self, node: TreeNode, val: int) -> TreeNode:
        if node is None:
            return TreeNode(val)
        if val < node.val:
            node.left = self._insert(node.left, val)
        elif val > node.val:
            node.right = self._insert(node.right, val)
        # equal: duplicate, ignored
        return node

    def search(self, val: int) -> TreeNode:
        return self._search(self.root, val)

    def _search(self, node: TreeNode, val: int) -> TreeNode:
        if node is None or node.val == val:
            return node
        if val < node.val:
            return self._search(node.left, val)
        return self._search(node.right, val)

    def delete(self, val: int) -> None:
        self.root = self._delete(self.root, val)

    def _delete(self, node: TreeNode, val: int) -> TreeNode:
        if node is None: return None
        if val < node.val:
            node.left = self._delete(node.left, val)
        elif val > node.val:
            node.right = self._delete(node.right, val)
        else: # found the match, now three cases
            if not node.left:
                return node.right
            if not node.right:
                return node.left
            #when both children exist
            # find inorder successor (smallest in the right subtree)
            successor = node.right
            while successor.left:
                successor = successor.left
            node.val = successor.val
            node.right = self._delete(node.right, successor.val)
        return node
\end{lstlisting}

\subsection{Validate a BST}

\begin{lstlisting}
# check if a given tree is actually BST
def is_valid_bst(node: TreeNode, low = float('-inf'), high = float('inf')) -> bool:
    if not node: return True

    if not (low < node.val < high): return False

    return is_valid_bst(node.left, low, node.val) and is_valid_bst(node.right, node.val, high)
\end{lstlisting}

\subsection{Lowest Common Ancestor (LCA)}

\begin{itemize}
\item The first version works for any binary tree (it assumes both \texttt{p} and \texttt{q} are present): $O(n)$.
\end{itemize}

\begin{lstlisting}
def lca(root:TreeNode, p:int, q:int) -> TreeNode:
    if root is None or root.val == p or root.val == q:
        return root

    left = lca(root.left, p, q)
    right = lca(root.right, p, q)

    if left and right: return root

    return left if left else right
\end{lstlisting}

For a BST we can use the ordering instead of searching both subtrees, which is $O(h)$:


\begin{lstlisting}
def lca_bst(root: TreeNode, p: int, q: int) -> TreeNode:
    while root:
        if p < root.val and q < root.val:
            root = root.left
        elif p > root.val and q > root.val:
            root = root.right
        else:
            return root   # p and q are split by root (or one equals root)
    return None
\end{lstlisting}

\section{Graphs}

\subsection{Basics}

\begin{itemize}
\item A graph is a set of nodes (vertices) connected via edges. Below, $V$ is the number of vertices and $E$ the number of edges.
\item No hierarchy, unlike trees.
\item Types of Graphs:
\begin{itemize}
\item Directed
\item Undirected
\item Weighted
\item Unweighted
\item Cyclic
\item Acyclic
\item Connected
\item Sparse vs Dense
\end{itemize}
\end{itemize}

\subsection{Representation}

\subsubsection{Adjacency List}

\begin{itemize}
\item A dictionary mapping each node to a list of  its neighbors. Efficient for sparse graphs.
\end{itemize}

\begin{lstlisting}
graph = {
    'A': ['B', 'C'],
    'B': ['A', 'D'],
    'C': ['A', 'D'],
    'D': ['B', 'C']
}
\end{lstlisting}

\begin{itemize}
\item For weighted graphs, we store \texttt{(\allowbreak{}neighbor,\allowbreak{} weight)} pairs.
\end{itemize}

\begin{lstlisting}
graph = {
    'A': [('B', 4), ('C', 1)],
    'B': [('A', 4), ('D', 2)],
    'C': [('A', 1), ('D', 5)],
    'D': [('B', 2), ('C', 5)]
}
\end{lstlisting}

\subsubsection{Adjacency Matrix}

\begin{itemize}
\item A 2d array, where \texttt{matrix[\allowbreak{}i][\allowbreak{}j] =\allowbreak{} 1} if there is an edge between \texttt{i} and \texttt{j}.
\end{itemize}

\begin{lstlisting}
# 4 nodes: 0, 1, 2, 3
matrix = [
    [0, 1, 1, 0],
    [1, 0, 0, 1],
    [1, 0, 0, 1],
    [0, 1, 1, 0]
]
\end{lstlisting}

\subsubsection{Edge List}

\begin{itemize}
\item Just a flat list of edges. (rarely used for traversal but common way of input in Qs)
\end{itemize}

\begin{lstlisting}
edges = [('A', 'B'), ('A', 'C'), ('B', 'D'), ('C', 'D')]
\end{lstlisting}

\sbox{\tblbox}{\begingroup\footnotesize\renewcommand{\arraystretch}{1.15}%
\setlength{\tabcolsep}{2.5pt}%
\setlength{\tw}{\dimexpr\columnwidth-4\tabcolsep\relax}%
\begin{tabular}{@{}>{\raggedright\arraybackslash}p{0.4618\tw}>{\raggedright\arraybackslash}p{0.2649\tw}>{\raggedright\arraybackslash}p{0.2733\tw}@{}}
\toprule
 & \textbf{Adjacency list} & \textbf{Adjacency matrix} \\
\midrule
Space & $O(V+E)$ & $O(V^2)$ \\
Is there an edge u–v? & $O(\mathrm{deg}(u))$ & $O(1)$ \\
Iterate over neighbours of u & $O(\mathrm{deg}(u))$ & $O(V)$ \\
Best for & sparse graphs & dense graphs \\
\bottomrule
\end{tabular}\endgroup}

\needspace{\dimexpr\ht\tblbox+\dp\tblbox+4.1\baselineskip\relax}

\noindent\textbf{Adjacency list vs. adjacency matrix}\par\nopagebreak\vspace{1pt}

\usebox{\tblbox}\smallskip

\subsubsection{Building an adjacency list from edge list}

\begin{lstlisting}
from collections import defaultdict
from typing import List

def build_adj_list(edges: List, directed: bool = False) -> dict:
    graph = defaultdict(list)
    for u, v in edges:
        graph[u].append(v)
        if directed:
            graph[v]                 # touch v, so that sink nodes also become keys
        else:
            graph[v].append(u)
    return graph

# the same with a plain dict (no import needed)
def build_adj_list_plain(edges: List, directed: bool = False) -> dict:
    graph = {}
    for u, v in edges:
        graph.setdefault(u, []).append(v)
        if directed:
            graph.setdefault(v, [])
        else:
            graph.setdefault(v, []).append(u)
    return graph
\end{lstlisting}

\subsection{Traversals}

\begin{itemize}
\item DFS and BFS both run in $O(V+E)$ time on an adjacency list; \texttt{visited} prevents revisiting nodes (and looping forever on cycles).
\item Recursive DFS is limited by Python's recursion limit (about 1000 frames by default); use the iterative version for deep graphs.
\item The cycle detection, topological sort and traversal code below looks up \texttt{graph[\allowbreak{}node]} for every node, so every node must be a key of \texttt{graph} (this is why the builder above also adds sink nodes).
\end{itemize}

\subsubsection{DFS (Recursive)}

\begin{lstlisting}
def dfs(graph, node, visited=None, result=None):
    if visited is None:
        visited, result = set(), []
    visited.add(node)
    result.append(node)
    for neighbor in graph[node]:
        if neighbor not in visited:
            dfs(graph, neighbor, visited, result)
    return result
\end{lstlisting}

\subsubsection{DFS (Iterative using Stack)}

\begin{itemize}
\item Visits the neighbours in the opposite order to the recursive version (the last one pushed is popped first).
\end{itemize}

\begin{lstlisting}
def dfs_iterative(graph, start):
    visited = set()
    stack = [start]
    result = []
    while stack:
        node = stack.pop()
        if node not in visited:
            visited.add(node)
            result.append(node)
            for neighbor in graph[node]:
                if neighbor not in visited:
                    stack.append(neighbor)
    return result
\end{lstlisting}

\subsubsection{BFS (always iterative, uses Queue)}

\begin{itemize}
\item BFS visits nodes in order of distance from \texttt{start}, so it finds shortest paths in unweighted graphs.
\end{itemize}

\begin{lstlisting}
from collections import deque

def bfs(graph, start):
    visited = {start}
    q = deque([start])
    result = []
    while q:
        node = q.popleft()
        result.append(node)
        for neighbor in graph[node]:
            if neighbor not in visited:
                visited.add(neighbor)
                q.append(neighbor)
    return result
\end{lstlisting}

\subsection{Common Algorithms}

\subsubsection{Detect a Cycle (Undirected Graph)}

\begin{itemize}
\item DFS; a visited neighbour that is not the parent means a cycle. (Parallel edges between the same two nodes are not detected by this version.)
\end{itemize}

\begin{lstlisting}
def detect_cycle_undirected(graph):
    visited = set()

    def dfs(node, parent):
        visited.add(node)
        for neighbor in graph[node]:
            if neighbor not in visited:
                if dfs(neighbor, node):
                    return True
            elif neighbor != parent:
                return True
        return False

    for node in graph:
        if node not in visited:
            if dfs(node, None): return True
    return False
\end{lstlisting}

\subsubsection{Detect a Cycle (Directed Graph)}

\begin{itemize}
\item Three colours: WHITE = unvisited, GRAY = on the current DFS path, BLACK = finished. An edge to a GRAY node closes a cycle.
\end{itemize}

\begin{lstlisting}
def detect_cycle_directed(graph):
    WHITE, GRAY, BLACK = 0, 1, 2
    color = {node: WHITE for node in graph}

    def dfs(node):
        color[node] = GRAY
        for neighbor in graph[node]:
            if color[neighbor] == GRAY:
                return True
            if color[neighbor] == WHITE and dfs(neighbor):
                return True
        color[node] = BLACK
        return False

    for node in graph:
        if color[node] == WHITE:
            if dfs(node): return True
    return False
\end{lstlisting}

\subsubsection{Topological Sort (DAGs Only)}

\begin{itemize}
\item An ordering in which, for every edge \texttt{u -> v}, \texttt{u} comes before \texttt{v}. DFS gives it as the reverse of the finishing order; $O(V+E)$. On a graph with a cycle the result is meaningless, so check for cycles first (or use Kahn's algorithm).
\end{itemize}

\begin{lstlisting}
def topological_sort(graph):
    visited = set()
    stack = []

    def dfs(node):
        visited.add(node)
        for neighbor in graph[node]:
            if neighbor not in visited:
                dfs(neighbor)
        stack.append(node) # add after all descendants are processed

    for node in graph:
        if node not in visited:
            dfs(node)

    return stack[::-1]
\end{lstlisting}

\subsubsection{Kahn's Algorithm (BFS alternative, also detects cycles)}

\begin{itemize}
\item Repeatedly remove nodes with in-degree 0. If some nodes are never removed, they lie on or behind a cycle.
\end{itemize}

\begin{lstlisting}
from collections import deque

def topological_sort_kahn(graph, nodes):
    in_degree = {n: 0 for n in nodes}
    for node in graph:
        for neighbor in graph[node]:
            in_degree[neighbor] += 1

    q = deque([n for n in nodes if in_degree[n] == 0])
    result = []

    while q:
        node = q.popleft()
        result.append(node)
        for neighbor in graph.get(node, []):
            in_degree[neighbor] -= 1
            if in_degree[neighbor] == 0:
                q.append(neighbor)

    if len(result) != len(nodes):
        return None   # cycle detected, no valid topological order
    return result
\end{lstlisting}

\subsubsection{Number of connected Components}

\begin{itemize}
\item For undirected graphs: every DFS started from an unvisited node explores exactly one component.
\end{itemize}

\begin{lstlisting}
def count_components(graph, nodes):
    visited = set()
    count = 0

    def dfs(node):
        visited.add(node)
        for neighbor in graph.get(node, []):
            if neighbor not in visited:
                dfs(neighbor)

    for node in nodes:
        if node not in visited:
            dfs(node)
            count += 1
    return count
\end{lstlisting}

\subsection{Shortest Paths}

\subsubsection{Dijkstra's Algorithm (only for non-negative weights)}

\begin{itemize}
\item Greedy: always settle the closest unsettled node, using a min-heap. $O((V+E) \log V)$.
\item Heap entries that became outdated are skipped when popped (\texttt{d > dist[\allowbreak{}node]}) instead of being removed. Unreachable nodes are simply absent from \texttt{dist}.
\end{itemize}

\begin{lstlisting}
import heapq

def dijkstra(graph, start):
    dist = {start: 0}
    heap = [(0, start)]

    while heap:
        d, node = heapq.heappop(heap)
        if d > dist.get(node, float('inf')):
            continue                  # stale entry, a shorter path was already found
        for neighbor, weight in graph.get(node, []):
            new_dist = d + weight
            if new_dist < dist.get(neighbor, float('inf')):
                dist[neighbor] = new_dist
                heapq.heappush(heap, (new_dist, neighbor))
    return dist
\end{lstlisting}

\subsubsection{Bellman-Ford Algorithm}

\begin{itemize}
\item Single-source shortest paths that \textbf{allows negative edge weights}. Relax every edge $V-1$ times; if some edge can still be relaxed afterwards, a negative cycle is reachable from the source. $O(VE)$.
\end{itemize}

\begin{lstlisting}
def bellman_ford(nodes, edges, start):
    """edges: list of (u, v, w). Returns the dist dict, or None on a negative cycle."""
    dist = {n: float('inf') for n in nodes}
    dist[start] = 0
    for _ in range(len(dist) - 1):
        updated = False
        for u, v, w in edges:
            if dist[u] + w < dist[v]:
                dist[v] = dist[u] + w
                updated = True
        if not updated:
            break                      # converged early
    for u, v, w in edges:
        if dist[u] + w < dist[v]:
            return None                # still improving -> negative cycle
    return dist
\end{lstlisting}

\subsubsection{SPFA - Shortest Path Faster Algorithm}

\begin{itemize}
\item Queue-based Bellman-Ford: only re-process a node when its distance has just improved. Fast on typical graphs, but still $O(VE)$ in the worst case. Also handles negative edges.
\item A node entering the queue $V$ times means a negative cycle.
\end{itemize}

\begin{lstlisting}
from collections import deque

def spfa(graph, nodes, start):
    """graph[u] = list of (v, w). Returns the dist dict, or None on a negative cycle."""
    dist = {n: float('inf') for n in nodes}
    dist[start] = 0
    q = deque([start])
    in_queue = {start}
    count = {n: 0 for n in nodes}      # how often a node was enqueued
    while q:
        u = q.popleft()
        in_queue.discard(u)
        for v, w in graph.get(u, []):
            if dist[u] + w < dist[v]:
                dist[v] = dist[u] + w
                if v not in in_queue:
                    count[v] += 1
                    if count[v] >= len(dist):
                        return None    # negative cycle
                    q.append(v)
                    in_queue.add(v)
    return dist
\end{lstlisting}

\subsubsection{Floyd-Warshall Algorithm}

\begin{itemize}
\item All-pairs shortest paths: after processing \texttt{k}, \texttt{dist[\allowbreak{}i][\allowbreak{}j]} is the shortest path from \texttt{i} to \texttt{j} that only uses intermediate nodes \texttt{0.\allowbreak{}.\allowbreak{}k}. The \texttt{k} loop must be the outermost one.
\item $O(V^3)$ time, $O(V^2)$ space; works with negative edges. A negative cycle exists if some \texttt{dist[\allowbreak{}i][\allowbreak{}i] < 0} at the end.
\end{itemize}

\begin{lstlisting}
def floyd_warshall(n, edges):
    """Nodes are 0..n-1, edges are (u, v, w). Returns the n x n distance matrix."""
    INF = float('inf')
    dist = [[INF] * n for _ in range(n)]
    for i in range(n):
        dist[i][i] = 0
    for u, v, w in edges:
        dist[u][v] = min(dist[u][v], w)
    for k in range(n):
        for i in range(n):
            for j in range(n):
                if dist[i][k] + dist[k][j] < dist[i][j]:
                    dist[i][j] = dist[i][k] + dist[k][j]
    return dist
\end{lstlisting}

\subsection{Minimum Spanning Tree}

\begin{itemize}
\item For a connected, weighted, undirected graph: $V-1$ edges that connect all vertices with the smallest possible total weight (no cycles).
\item Two classic greedy algorithms: Prim grows one tree from a start node by always taking the cheapest edge leaving it; Kruskal sorts all edges and adds each one that does not close a cycle.
\end{itemize}

\subsubsection{Prim's Algorithm}

\begin{itemize}
\item Heap of candidate edges. $O(E \log V)$. A good fit when the graph is given as an adjacency list.
\end{itemize}

\begin{lstlisting}
import heapq

def prim(graph, start):
    """graph[u] = list of (v, w), undirected (each edge stored in both directions).
    Returns (total_weight, mst_edges) for the component containing start."""
    visited = {start}
    heap = [(w, start, v) for v, w in graph[start]]
    heapq.heapify(heap)
    total, mst = 0, []
    while heap:
        w, u, v = heapq.heappop(heap)
        if v in visited:
            continue
        visited.add(v)
        total += w
        mst.append((u, v, w))
        for nxt, nw in graph[v]:
            if nxt not in visited:
                heapq.heappush(heap, (nw, v, nxt))
    return total, mst
\end{lstlisting}

\subsubsection{Kruskal's Algorithm}

\begin{itemize}
\item Sorting dominates: $O(E \log E)$. Cycle checks use \textbf{union-find} (disjoint set union), which is nearly $O(1)$ per operation with path compression and union by size. A good fit when the graph is given as an edge list.
\end{itemize}

\begin{lstlisting}
class DSU:
    def __init__(self, nodes):
        self.parent = {n: n for n in nodes}
        self.size = {n: 1 for n in nodes}

    def find(self, x):
        while self.parent[x] != x:
            self.parent[x] = self.parent[self.parent[x]]   # path halving
            x = self.parent[x]
        return x

    def union(self, a, b):
        ra, rb = self.find(a), self.find(b)
        if ra == rb:
            return False               # already connected, would form a cycle
        if self.size[ra] < self.size[rb]:
            ra, rb = rb, ra
        self.parent[rb] = ra
        self.size[ra] += self.size[rb]
        return True

def kruskal(nodes, edges):
    """edges: list of (u, v, w). Returns (total_weight, mst_edges); a forest if disconnected."""
    dsu = DSU(nodes)
    total, mst = 0, []
    for u, v, w in sorted(edges, key=lambda e: e[2]):
        if dsu.union(u, v):
            total += w
            mst.append((u, v, w))
    return total, mst
\end{lstlisting}

\subsection{Minimum Height Trees}

\begin{itemize}
\item Given a tree (connected, acyclic, undirected) with $n$ nodes, choose the root that minimises the height of the rooted tree. There are at most two such roots: the centre(s) of the tree.
\item Idea: repeatedly peel off all leaves, like Kahn's algorithm working from the outside in. The one or two nodes left at the end are the answer. $O(n)$.
\end{itemize}

\begin{lstlisting}
def find_min_height_trees(n, edges):
    if n <= 2:
        return list(range(n))
    adj = [set() for _ in range(n)]
    for u, v in edges:
        adj[u].add(v)
        adj[v].add(u)
    leaves = [i for i in range(n) if len(adj[i]) == 1]
    remaining = n
    while remaining > 2:
        remaining -= len(leaves)
        new_leaves = []
        for leaf in leaves:
            nb = adj[leaf].pop()
            adj[nb].remove(leaf)
            if len(adj[nb]) == 1:
                new_leaves.append(nb)
        leaves = new_leaves
    return leaves
\end{lstlisting}

\section{Bit Manipulation}

\subsection{Basics}

\subsubsection{Basic Operators}

\begin{lstlisting}
a = 5           # 0101
b = 3           # 0011

a & b           # 1,  0001, and
a | b           # 7,  0111, or
a ^ b           # 6,  0110, xor

~a              # -6, flips all bits; in Python, ~a is defined as -(a+1)

a << 1          # 10, 1010, multiply by 2
a >> 1          # 2,  0010, divide by 2 (floor)
\end{lstlisting}

\subsubsection{XOR Properties}

\begin{itemize}
\item XOR is commutative and associative, so the order does not matter and equal pairs cancel out.
\end{itemize}

\begin{lstlisting}
a ^ a           # 0, xor anything with itself is 0
a ^ 0           # a, xor with 0 preserves the value
a ^ a ^ b       # b, xor-ing twice with the same value cancels out
\end{lstlisting}

\subsubsection{NOT Properties}

\begin{itemize}
\item Python integers have arbitrary precision (there is no fixed width), so \texttt{\textasciitilde{}a} cannot flip \emph{all} bits; it is defined as ${\sim}a = -(a + 1)$ instead.
\item If we want a fixed-width bit-flipped version, we can mask the result:
\end{itemize}

\begin{lstlisting}
def not_8bit(a):
    return ~a & 0xFF # 0xFF = 11111111, keeps only last 8 digits
\end{lstlisting}

\subsubsection{Shifting}

\begin{lstlisting}
1 << n          # 2^n
x >> n          # floor(x / 2^n)
\end{lstlisting}

\subsection{Common Algorithms}

\subsubsection{Check if the i-th bit is set}

\begin{lstlisting}
def is_bit_set(n,i):
    return (n>>i) & 1 == 1
\end{lstlisting}

\subsubsection{Set the i-th bit}

\begin{lstlisting}
def set_bit(n,i):
    return n | (1 << i)
\end{lstlisting}

\subsubsection{Clear the i-th bit}

\begin{lstlisting}
def clear_bit(n,i):
    return n & ~(1 << i)
\end{lstlisting}

\subsubsection{Toggle the i-th bit}

\begin{lstlisting}
def toggle_bit(n,i):
    return n ^ (1 << i)
\end{lstlisting}

\subsubsection{Check if a number is a power of 2}

\begin{lstlisting}
def is_power_of_two(n):
    return n > 0 and (n & (n - 1)) == 0
\end{lstlisting}

\subsubsection{Count set bits}

\begin{itemize}
\item \texttt{n \& (\allowbreak{}n - 1)} clears the lowest set bit, so the loop runs once per set bit. (This is also why the power-of-two test above works.)
\end{itemize}

\begin{lstlisting}
def count_bits(n):
    count = 0
    while n:
        n &= n - 1
        count += 1
    return count
\end{lstlisting}

\subsubsection{XOR based pattern problems}

\begin{lstlisting}
# every number appears twice except one: find the single one
from typing import List

def findSingle(nums: List[int]) -> int:
    ans = 0
    for x in nums:
        ans ^= x
    return ans
\end{lstlisting}

\begin{lstlisting}
# find the missing number in a range 0-n
from typing import List

def findMissing(nums: List[int]) -> int:
    result = len(nums)
    for i, num in enumerate(nums):
        result ^= i^num
    return result
\end{lstlisting}

\begin{lstlisting}
# swapping two numbers without a temporary (in Python, a, b = b, a is the idiomatic way)
a, b = 5, 3
a = a ^ b
b = a ^ b # (a ^ b) ^ b = a
a = a ^ b # (a ^ b) ^ ((a ^ b) ^ b) = b
\end{lstlisting}

\begin{lstlisting}
# all subsets
from typing import List

def allSubsets(nums: List[int]) -> List[List[int]]:
    n = len(nums)
    subsets = []

    for mask in range(1 << n): # 2^n choices
        subset = [nums[i] for i in range(n) if mask & (1 << i)]
        subsets.append(subset)

    return subsets
\end{lstlisting}
\end{multicols}
\end{document}
